\documentclass[11pt]{article}
\usepackage[margin=1in]{geometry}
\usepackage[T1]{fontenc}
\usepackage{lmodern}
\usepackage{microtype}
\usepackage{url}

\begin{document}

\section{Introduction}

Automated guided vehicles (AGVs) are a central part of material handling in
smart manufacturing. They move workpieces between storage, processing, and
delivery locations, allowing production systems to respond flexibly to
changing orders and workloads. As factories grow, however, coordinating a
fleet becomes a joint assignment, routing, and traffic-management problem.
Vehicles may need the same narrow aisle, workstation approach, or transport
resource at overlapping times. A schedule must therefore assign jobs and
produce routes and timings that avoid collisions while meeting operational
targets such as completion time, delivery due dates, and balanced vehicle
utilization. These requirements become more demanding in multi-floor
facilities, where lifts connect otherwise separate floor networks. A lift is
a shared resource with finite capacity and non-negligible travel time: a
conflict at one lift can delay routes on several floors and propagate through
the entire schedule.

Many AGV scheduling methods centralize assignment and traffic planning. This
provides a system-wide view, but it also concentrates computation and
reservation state at one server. As the fleet and route set grow, that server
must reason about an increasing number of node and track occupancy windows,
including interactions among vehicles operating in different areas. The
planning workload and shared state can consequently become a bottleneck.
This challenge is especially relevant to multi-floor factories: a central
planner must coordinate local floor traffic as well as cross-floor lift
movements. A practical architecture should retain enough central coordination
to prevent conflicting lift use while distributing the detailed routing and
traffic state across the factory.

Guo et al. present a prioritized planning method for large-scale multiple-AGV
scheduling in smart manufacturing \cite{guo2024}. Their approach assigns
transport tasks and plans AGV routes sequentially according to priority,
resolving conflicts with respect to plans already made for higher-priority
vehicles. This provides the centralized baseline for the present work. In
that baseline, assignment and conflict-free routing are handled by a single
planning process. The paper and its accompanying walkthrough material provide
the algorithmic notation and worked examples used here
\cite{guo2024,walkthrough}. The repository contains a copy of the paper as
\texttt{Prioritized\_Planning\_for\_Large-Scale\_Multiple-AGV\_Scheduling\_Problem\_in\_Smart\_Manufacturing.pdf}
and the walkthrough presentation as
\texttt{Prioritized\_Planning\_AGV\_Walkthrough\_with\_Extension.pptx}.

This paper extends that prioritized-planning baseline with a hybrid
central-plus-edge architecture for a multi-floor factory. A central server
performs job assignment and divides each AGV route into floor-level legs. It
also books the lift trips, which remain shared across floors and therefore
require system-wide coordination. Each floor has an edge server that plans
its local route legs and resolves node and track conflicts using the
prioritized conflict-free routing procedure. The decomposition limits each
floor server's detailed reservation state to its own floor, while the central
server retains responsibility for assignments and inter-floor transitions.
To give vehicles a safe place to wait near a local conflict or a busy lift,
the layout adds a lobby node at each lift landing. A floor server can delay a
leg at its entry lobby, and the central server can schedule a lift departure
after the lift resource becomes available. These lobbies provide waiting
points near the conflict rather than requiring every delay to be absorbed at
the factory's single parking node.

The implementation is evaluated against two centralized configurations: the
paper-style planner, and a centralized planner with the same lift lobbies as
the hybrid system. The saved comparison uses 100 jobs and 20 AGVs over 10
random seeds; each generated plan is checked by an independent verifier. In
the reported mean results, the hybrid configuration reduces last-delivery
makespan by 65\%, total tardiness by 72\%, and usage-time imbalance by 80\%
relative to the paper-style centralized configuration. Its busiest-server
planning CPU time is 93\% lower. These figures describe the repository's
modeled factory, parameters, and tested seeds; they should be read as
experimental results for this setup, not as general guarantees for all
multi-floor factories. The comparison also includes a fleet-size sweep over
4, 8, 12, 16, and 20 AGVs, providing a view of how the measured scheduling
metrics vary with fleet size \cite{results}.

The remainder of the paper is organized as follows. Section~2 defines the
factory graph, job-assignment rules, route windows, conflict constraints, and
performance measures. Section~3 describes the centralized baseline and the
hybrid central and floor-server procedures, including lift reservations and
lobby waiting. Section~4 presents the experimental setup, verification
method, and results over the ten seeds and fleet-size sweep. Section~5
discusses the results, modeling assumptions, and limitations. Section~6
concludes the paper and identifies directions for extending the architecture.

\section*{Contributions}
\begin{itemize}
  \item Reimplementation of the ICARCV 2024 prioritized-planning baseline,
        including fixes to the original demo's documented bugs.
  \item A hybrid architecture that separates central job assignment and lift
        booking from per-floor conflict-free route planning.
  \item Lift-lobby waiting points that allow delays to be absorbed near
        floor-level conflicts and lift resources.
  \item Independent verification of completed plans for node and track
        conflicts, edge travel times, dwell windows, route endpoints, and
        exactly-once job assignment.
  \item Comparative results over ten random seeds and a fleet-size sweep.
\end{itemize}

\begin{thebibliography}{9}

\bibitem{guo2024}
Y.~Guo, J.~Yao, J.~Li, R.~Su, and K.~V.~Ling,
``Prioritized Planning for Large-Scale Multiple-AGV Scheduling Problem in
Smart Manufacturing,'' in \emph{Proceedings of the 18th International
Conference on Control, Automation, Robotics and Vision (ICARCV)}, 2024,
pp.~19--24, doi:~\url{https://doi.org/10.1109/ICARCV63323.2024.10821683}.

\bibitem{walkthrough}
\emph{Prioritized Planning AGV Walkthrough with Extension}, project
walkthrough presentation, 51 slides, repository file
\texttt{Prioritized\_Planning\_AGV\_Walkthrough\_with\_Extension.pptx}.
The repository also includes a local copy of the ICARCV paper PDF under the
filename cited in the Introduction.

\bibitem{results}
Project repository, \texttt{agv\_code/agv\_code/results\_hybrid/summary.txt},
saved comparison summary for 100 jobs, 20 AGVs, and seeds 1--10.

\end{thebibliography}

\end{document}
